\subsection{Naturalité de l’échelle et formes canoniques opératorielles}
\label{rev5:naturalidad-areal}

La longueur engendrée permet de préciser la compatibilité entre la
réalisation positive, l’opérateur d’aire et l’information sectorielle.
La représentation marquée \(\mathcal F(K,\xi)\) conserve,
parmi les données pertinentes de \(\xi\), la résolution
\(\Omega\), ses projecteurs, le retour transporté et la section
\(L_*\). Les coordonnées régionales engendrées, l’action et les canaux
constitutifs maintiennent leur généalogie antérieure. L’inverse marqué
restitue ainsi tous les arguments de
\eqref{rev5:longitud-generada} et de \eqref{rev5:G-area}.
En particulier, le lecteur géométrique déjà construit de
\(\alpha^{\mathrm{geom}}\) détermine les expressions
\[
 L_\alpha^{\mathrm{geom}}
 =L_*(\alpha^{\mathrm{geom}})^{16}\frac{5759}{23040},
 \qquad
 G^{\mathrm{geom}}
 =\frac{(c^{\mathrm{geom}})^3
 (L_\alpha^{\mathrm{geom}})^2}{\hbar_{\rm ret}^{\mathrm{geom}}}.
\]
Les lecteurs de vitesse et d’action s’obtiennent par la même
composition avec l’inverse marqué. Ces expressions conservent
les données du retour et la section radiale de référence.

\begin{corollary}[Conservation par changements de coordinatisation et raffinement hérité]
\label{rev5:conservacion}
Les changements positifs de coordinatisation conservant le point marqué et les raffinements restituant le registre initial conservent \(L_\alpha\), \(G\), \(\gamma_{\rm HMT}\) et \(\widehat{\mathcal A}=L_\alpha^2a_\partial N\) lorsqu'ils transportent également les données sectorielles et dimensionnelles déclarées. Le spectre normalisé \(\widehat{\mathcal A}/L_\alpha^2=a_\partial N\), ses multiplicités et l'état \(\rho_A(\Delta)\) demeurent identiques.
\end{corollary}
\begin{proof}
Le théorème~\ref{thm:compat-algebra-lectores} transporte chaque lecteur
par composition avec l’inverse. Les mutations du même point
reconstruisent le même registre. Dans le raffinement hérité, la somme
des séparations subdivisées reconstruit chaque coordonnée initiale.
La tensorisation de l’inscription conserve \(r_\Omega\),
et ses autres arguments restent fixés par les marques. Par
substitution, la longueur, le coefficient gravitationnel et
l’opérateur d’aire sont conservés. Les projecteurs d’occupation conservent le
polynôme de multiplicités \((1+z+z^2)^{36}\) et l’expression
normalisée de l’état.
\end{proof}

Le raffinement hérité considéré ici conserve le même espace
de trente-six liens de frontière. Une extension qui introduit
un espace physique de frontière différent requiert des opérateurs d’entrelacement
spécifiés pour \(N_{90}\) et \(N_{120}\), ainsi que le
transport des données sectorielles déclarées.

Soit maintenant \(P\) un polytope de rang un de ce manuscrit, avec une subdivision orientée \(P=\bigcup_\sigma P_\sigma\) du même support, de multiplicité un et à faces compatibles. On fixe la même fibre sectorielle \(\mathscr H_A\) sur toutes ses cellules. La forme de degré maximal à valeurs opératorielles et de dimension d'aire est définie par
\[
 \boldsymbol\Omega^{\mathcal A}_P
 =\widehat{\mathcal A}\otimes\Omega_P.
\]
Ici \(\widehat{\mathcal A}\) est constant par rapport aux
coordonnées différentielles du polytope : il provient de l’état HMT
fixé, tandis que \(\Omega_P\) décrit son support géométrique.
Le coefficient \(\widehat{\mathcal A}\) est semi-défini positif.
Sa composante dans le secteur du vide \(N=0\) est nulle. Sur un
espace propre d’occupation \(n>0\), diviser par
\(nL_\alpha^2a_\partial\) récupère la forme canonique
\(\Omega_P\) multipliée par l’identité de cet espace propre.

\begin{proposition}[Additivité aréale et résidus de frontière]
\label{rev5:forma-areal}
Dans la trivialisation sectorielle commune, on a
\begin{equation}
 \boldsymbol\Omega^{\mathcal A}_P
 =\sum_\sigma\widehat{\mathcal A}\otimes\Omega_{P_\sigma},
 \qquad
 \operatorname{Res}_F\boldsymbol\Omega^{\mathcal A}_P
 =\widehat{\mathcal A}\otimes\Omega_F.
\end{equation}
Les identités persistent sous subdivisions successives et résidus itérés. Pour un état fixé \(\rho\), prendre la trace donne \(\Tr(\rho\widehat{\mathcal A})\Omega_P\) et conserve les mêmes identités scalaires.
\end{proposition}
\begin{proof}
L’espace d’opérateurs de \(\mathscr H_A\) est de dimension
finie. Dans toute base, l’assertion est l’identité
d’additivité des formes canoniques multipliée par chaque coefficient
constant de \(\widehat{\mathcal A}\). Le résidu et la trace
sont linéaires. L’annulation des facettes intérieures conserve leurs
orientations et le même coefficient opératoriel des deux côtés.
L’itération applique ces opérations à chaque subdivision et à chaque
drapeau de faces.
\end{proof}

Si différentes cellules portent des coefficients opératoriels réguliers \(A_\sigma\), le défaut d'une paire sur sa face commune est \((A_\sigma|_F-A_\tau|_F)\otimes\Omega_F\). L'égalité des restrictions opératorielles à la frontière fournit le recollement, sous les conditions de pôles simples du théorème~\ref{thm:compat-pegado}. La somme globale conserve en outre toutes les contributions sur chaque diviseur ambiant et le contrôle des pôles extérieurs. La normalisation dimensionnelle, l'annulation des résidus et la couverture du support ont ainsi des critères séparés et compatibles.

L’identité modulaire finie conserve aussi une expression directe
en aires physiques \(\mathcal A_m=L_\alpha^2a_\partial N_m\) :
\begin{equation}
 s(\sigma)-s(\rho_0)
 =\sum_{m=90,120}\frac{m\Delta_0}{L_\alpha^2a_\partial}
   \Tr[(\sigma-\rho_0)\mathcal A_m]
  -D(\sigma\Vert\rho_0).
 \label{rev5:modular-fisica}
\end{equation}
C’est \eqref{rev4:ae:ley-finita-formula} après insertion de la
longueur construite. Le facteur angulaire \(a_0\), contenu dans
\(a_\partial=\gamma_{\rm HMT}a_0\), demeure intégral.
Si l’on compare des sections dimensionnelles avec \(L'_\alpha=sL_\alpha\)
et si l’on conserve le même état et les mêmes données adimensionnelles,
les aires sont multipliées par \(s^2\) et leurs coefficients modulaires
par \(s^{-2}\). L’entropie et l’entropie relative demeurent
égales. Cette comparaison de sections ne constitue pas une autre solution
pour les mêmes données longitudinales fixées.

Pour un horizon sphérique de rayon \(R_h=L_\alpha\zeta_h\),
l’aire géométrique vaut \(4\pi_{\rm HMT}L_\alpha^2\zeta_h^2\).
Son identification à l’espérance de l’opérateur sectoriel, si elle est
requise pour une même coupe, exige la condition explicite
\[
 a_\partial\Tr(\rho N)=4\pi_{\rm HMT}\zeta_h^2.
\]
La géométrie de l'horizon et l'état de frontière doivent satisfaire cette relation avec leurs propres constructions. Le facteur \(\zeta_h\) caractérise cet horizon ; l'échelle longitudinale du couplage reste \(L_\alpha\).
